\subsection{Part A – Knowledge check}

\indent 

\noindent\textbf{What is the difference between Project Management (PM) and Change Management (CM) in IT projects?}


PM: Focuses on delivering the technical side of the project — scope, budget, schedule, quality, and resources.

CM: Focuses on the people side — ensuring employees adopt and use the new systems/processes.


\bigskip

\noindent\textbf{Why do organizations need change management even if the project is delivered on time and within budget?}



A project can “succeed” technically but fail in value realization if people resist or ignore the change.

CM drives adoption, utilization, and proficiency → without it, ROI is not achieved, and investments are wasted.


\bigskip

\noindent\textbf{Briefly explain one change management model (McKinsey 7S, ADKAR, or Kotter’s 8-Step) and when it would be most useful.}



\noindent\textbf{McKinsey 7S Model}

A holistic organizational framework with 7 elements: Strategy, Structure, Systems, Shared Values, Skills, Style, and Staff.

Helps diagnose misalignments that block change.

When useful: Large-scale IT or organizational transformations where both hard (strategy, systems) and soft (skills, culture, leadership style) factors need alignment.

\bigskip

\noindent\textbf{ADKAR Model}

Developed by Prosci, focusing on individual adoption of change. 

Stages: Awareness, Desire, Knowledge, Ability, Reinforcement.

Ensures people move through the process of accepting and sustaining change.

When useful: IT rollouts where end-user adoption is critical and resistance is expected.

\bigskip

\noindent\textbf{Kotter’s 8-Step Model}

A top-down, leadership-driven model that emphasizes building momentum for change.

Steps include: Create urgency, Build coalition, Develop vision, Communicate vision, Empower action, Generate wins, Consolidate gains, Anchor in culture.

When useful: Organization-wide cultural or strategic changes where leadership sponsorship and long-term embedding are key


\subsection{Part B – Case Study Discussion}

BetaBank, a mid-sized retail bank, is implementing a new Customer Relationship Management (CRM) system to improve customer service and enable data-driven marketing. The background of this project is as follows:

\begin{itemize}
    \item The IT team is excited about the system’s AI-driven analytics.
    \item The Project Managers are confident the system will go live on time.
    \item Branch staff worry the CRM is too complex and will slow down customer interactions.
    \item Marketing staff don’t understand how the new system differs from their current spreadsheets.
    \item After the pilot phase, adoption rates are only 30\%, with many employees reverting to old tools.
    \item The CEO is frustrated: “We’ve spent \$5 million, but the results are not visible.”
\end{itemize}

\bigskip

\noindent Based on the above case study, discuss as a group the following:

\noindent\textbf{At which stages of the change curve might the branch staff and marketing team be?}



Branch staff: Likely in Denial (“the old system works fine”) or Anger (“this will slow us down”).

Marketing staff: Likely in Shock/Confusion or Bargaining (“our spreadsheets are simpler; can we keep using them?”).


\bigskip

\noindent\textbf{What psychological barriers are visible in this case?}



Fear of loss (time, efficiency, comfort with old tools).

Uncertainty about how CRM benefits them directly.

Lack of trust (staff don’t see leadership addressing concerns).

Low motivation (they don’t see personal value in switching).


\bigskip

\noindent\textbf{Which change management model would you apply here? Why?}



ADKAR is most suitable: Staff lack Awareness (why CRM is better), have little Desire (don’t see personal benefit), and limited Knowledge/Ability (training not effective).

ADKAR directly addresses individual adoption barriers step by step.


\bigskip

\noindent\textbf{What are the potential risks to business outcomes if the CRM adoption continues to fail?}



Business risks: Poor customer service, missed cross-selling opportunities, wasted \$5M investment, employee frustration, reputational damage.

IT risks: Low adoption rates, shadow systems (spreadsheets), data inconsistency, governance failures.


\bigskip

\noindent\textbf{Suggest three specific interventions the Change Management team could take in the next 3 months to improve adoption.}



Targeted Communication \& Awareness Campaign: Explain benefits in employee language (“CRM reduces double data entry, saves time with customers”).

Tailored Training \& Support: Hands-on workshops, branch-level “CRM champions,” and peer-to-peer learning instead of generic sessions.

Reinforcement \& Incentives: Recognize and reward teams actively using CRM; set KPIs for usage; share success stories from early adopters
